\section{AHSME 1980}
        \begin{enumerate}
        	\item (\textit{AHSME 1980}) The largest whole number such that seven times the number is less than $100$ is


 $\text{(A)} \ 12 \qquad \text{(B)} \ 13 \qquad \text{(C)} \ 14 \qquad \text{(D)} \ 15 \qquad \text{(E)} \ 16$

 \textbf{Answer: }C



\textbf{Solution}

Let $x$ be a whole number such that $7x < 100$. Dividing by $7$ yields $x < \frac{100}{7} =  14\dfrac{2}{7}$, so $x = \boxed{(\textbf{C})\ 14}$ is the maximum possible value.

 


	\item (\textit{AHSME 1980}) The degree of $(x^2+1)^4 (x^3+1)^3$ as a polynomial in $x$ is


 $\text{(A)} \ 5 \qquad \text{(B)} \ 7 \qquad \text{(C)} \ 12 \qquad \text{(D)} \ 17 \qquad \text{(E)} \ 72$

 \textbf{Answer: }D



\textbf{Solution 1}

Expanding each factor yields $(x^{8}+...)(x^{9}+...)$. The degree of the first factor is $8$, while the degree of the second factor is $9$. Therefore, the degree of the polynomial is $8 + 9 = \boxed{(\textbf{D})\ 17}$




\textbf{Solution 2}

Let $\deg{P(x)}$ be the degree of a polynomial in $x$.  Recall that for any polynomials $P(x)$ and $Q(x)$, and nonnegative integer $n$, $\deg{(P(x)^n)} = n \deg{P(x)}$ and $\deg{(P(x)Q(x))} = \deg {P(x)} + \deg {Q(x)}$.

 So $\deg{((x^2+1)^4(x^3+1)^3)}$ $= \deg{((x^2+1)^4)} + \deg{((x^3+1)^3)}$ $= 4 \deg{(x^2+1)} + 3 \deg{(x^3+1)}$ $= 4 \cdot 2 + 3 \cdot 3$ $= \boxed{(\textbf{D})\ 17}$.

 -mihirb, edited by j314andrews

 


	\item (\textit{AHSME 1980}) If the ratio of $2x-y$ to $x+y$ is $\frac{2}{3}$, what is the ratio of $x$ to $y$?


 $\text{(A)} \ \frac{1}{5} \qquad \text{(B)} \ \frac{4}{5} \qquad \text{(C)} \ 1 \qquad \text{(D)} \ \frac{6}{5} \qquad \text{(E)} \ \frac{5}{4}$

 \textbf{Answer: }E



\textbf{Solution}

$\frac{2x-y}{x+y} = \frac{2}{3}$

 Cross multiplying yields  

 $6x-3y=2x+2y$
$4x=5y$
$\dfrac{x}{y}= \boxed{(\textbf{E}) \dfrac{5}{4}}$



	\item (\textit{AHSME 1980}) In the adjoining figure, $CDE$ is an equilateral triangle and $ABCD$ and $DEFG$ are squares. The measure of $\angle GDA$ is


 $\text{(A)} \ 90^\circ \qquad \text{(B)} \ 105^\circ \qquad \text{(C)} \ 120^\circ \qquad \text{(D)} \ 135^\circ \qquad \text{(E)} \ 150^\circ$

                    \begin{figure}[H]
                        \centering
                        \includegraphics[width=0.5\textwidth]{Problem/AHSME_I_1980_4_Problem_Figure_1.png}
                    \end{figure}
                    

 \textbf{Answer: }C



\textbf{Solution}

Since $ABCD$ and $DEFG$ are squares, $\angle ADC = \angle GDE = 90^\circ$.  Since $CDE$ is equilateral, $\angle CDE = 60^\circ$.  So $\angle GDA=360^\circ-90^\circ-60^\circ-90^\circ= \boxed{(\textbf{C})\ 120^\circ}$.

 


	\item (\textit{AHSME 1980}) If $AB$ and $CD$ are perpendicular diameters of circle $Q$, and $\measuredangle QPC = 60^\circ$, then the length of $PQ$ divided by the length of $AQ$ is 


 
                    \begin{figure}[H]
                        \centering
                        \includegraphics[width=0.5\textwidth]{Problem/AHSME_I_1980_5_Problem_Figure_1.png}
                    \end{figure}
                    
$\text{(A)} \ \frac{\sqrt{3}}{2} \qquad \text{(B)} \ \frac{\sqrt{3}}{3} \qquad \text{(C)} \ \frac{\sqrt{2}}{2} \qquad \text{(D)} \ \frac12 \qquad \text{(E)} \ \frac23$

 \textbf{Answer: }B



\textbf{Solution}

Notice that $\angle PCQ=30^\circ$, and therefore $\triangle PCQ$ is a $30^\circ-60^\circ-90^\circ$ right triangle. Let $x = PQ$. Then $CQ=AQ=x\sqrt{3}$ and $\frac{PQ}{AQ}=\frac{x}{x\sqrt{3}}=\boxed{(\textbf{B})\ \frac{\sqrt{3}}{3}}$.

 


	\item (\textit{AHSME 1980}) A positive number $x$ satisfies the inequality $\sqrt{x} < 2x$ if and only if


 $\text{(A)} \ x > \frac{1}{4} \qquad \text{(B)} \ x > 2 \qquad \text{(C)} \ x > 4 \qquad \text{(D)} \ x < \frac{1}{4}\qquad \text{(E)} \ x < 4$

 \textbf{Answer: }A



\textbf{Solution 1}

Since both sides of this inequality are positive, we may square both sides to get $x < 4x^2$.  Since $x$ is positive, we can divide both sides by $x$ to get $1 < 4x$.  Dividing both sides by $4$ yields $\frac{1}{4} < x$, or $\boxed{(\textbf{A})\ x > \frac{1}{4}}$.

 



\textbf{Solution 2 (Graphing)}

Graphing both $\sqrt{x}$ and $2x$ in Quadrant I of the coordinate plane, we see that the two curves intersect at $\left(\frac{1}{4}, \frac{1}{2}\right)$, and $2x > \sqrt{x}$ whenever $\boxed{(\textbf{A})\ x > \frac{1}{4}}$.

 


	\item (\textit{AHSME 1980}) Sides $AB,BC,CD$ and $DA$ of convex polygon $ABCD$ have lengths $3$, $4$, $12$, and $13$, respectively, and $\measuredangle CBA$ is a right angle. The area of the quadrilateral is


 
                    \begin{figure}[H]
                        \centering
                        \includegraphics[width=0.5\textwidth]{Problem/AHSME_I_1980_7_Problem_Figure_1.png}
                    \end{figure}
                    
$\text{(A)} \ 32 \qquad \text{(B)} \ 36 \qquad \text{(C)} \ 39 \qquad \text{(D)} \ 42 \qquad \text{(E)} \ 48$

 \textbf{Answer: }B



\textbf{Solution}

Draw $\overline{AC}$. By the Pythagorean Theorem, $AC = \sqrt{AB^2 + BC^2} = \sqrt{3^2+4^2} = 5$.  Notice that $AD^2 = AC^2 + CD^2 = 169$, so by the converse of the Pythagorean Theorem, $\triangle ACD$ is right.  Thus the area of $ABCD$ is $\frac{3\cdot4}{2}+\frac{5\cdot12}{2}=\boxed{(\textbf{B})\ 36}$.

 


	\item (\textit{AHSME 1980}) How many pairs $(a,b)$ of non-zero real numbers satisfy the equation


 
 \[\frac{1}{a} + \frac{1}{b} = \frac{1}{a+b}\]
$\text{(A)} \ \text{none} \qquad \text{(B)} \ 1 \qquad \text{(C)} \ 2 \qquad \text{(D)} \ \text{one pair for each} ~b \neq 0$
$\text{(E)} \ \text{two pairs for each} ~b \neq 0$

 \textbf{Answer: }A



\textbf{Solution}

Adding the two fractions on the left side yields $\frac{a+b}{ab}=\frac{1}{a+b}$.  

 Cross-multiplying yields $a^2+2ab+b^2=ab$.  

 Subtracting $ab$ from both sides yields $a^2+ab+b^2=0$.

 Suppose $b$ is constant and $a$ is variable.  Then the discriminant of this quadratic equation is $b^2 - 4 \cdot 1 \cdot b^2 = -3b^2$, which is negative if $b \neq 0$.  Therefore, for each $b \neq 0$, this equation has no real solutions, and the answer is $\boxed{\text{(\textbf{A})\ none}}$.

 


	\item (\textit{AHSME 1980}) A man walks $x$ miles due west, turns $150^\circ$ to his left and walks $3$ miles in the new direction. If he finishes a a point $\sqrt{3}$ from his starting point, then $x$ is


 $\text{(A)} \ \sqrt 3 \qquad \text{(B)} \ 2\sqrt{5} \qquad \text{(C)} \ \frac 32 \qquad \text{(D)} \ 3 \qquad \text{(E)} \ \text{not uniquely determined}$

 \textbf{Answer: }E



\textbf{Solution 1}


                    \begin{figure}[H]
                        \centering
                        \includegraphics[width=0.5\textwidth]{Solution/AHSME_I_1980_9_Solution_Figure_1.png}
                    \end{figure}
                    

 Let $\ell$ be the line he was initially walking on, $T$ be his turning point, $F$ be his finishing point, and $H$ be the foot of the perpendicular from $F$ to $\ell$.  Since he turned $150^\circ$,  $\angle HTF = 30^{\circ}$.  Therefore, $HF = \frac{3}{2}$ miles.  Since $F$ is $\sqrt{3}$ miles from his starting point, his starting point lies on a circle of radius $\sqrt{3}$ miles centered at $F$.  Since $\frac{3}{2} < \sqrt{3} < 3$, this circle must intersect $\ell$ once between $T$ and $H$ and once east of $H$.  Let $S_1$ and $S_2$ be these two intersections, respectively.  Both $S_1$ and $S_2$ are possible starting points, and clearly $TS_1 \neq TS_2$, so $x$ is $\fbox{(\textbf{E}) not uniquely determined}$.

 



\textbf{Solution 2}


                    \begin{figure}[H]
                        \centering
                        \includegraphics[width=0.5\textwidth]{Solution/AHSME_I_1980_9_Solution_Figure_2.png}
                    \end{figure}
                    

 
Let $S$ be his starting point, $T$ be the point where he turns, and $F$ be his finishing point.  Since he turned $150^{\circ}$ at $T$, $\angle STF = 30^{\circ}$.  By the Law of Cosines, $FS^2 = FT^2 + ST^2 - 2 \cdot FT \cdot ST \cos \angle STF$. That is, $\left(\sqrt{3}\right)^2 = 3^2 + x^2 - 2 \cdot 3 \cdot x \cdot \cos 30^{\circ}$.  Combining all terms on one side yields $x^2 - 3x\sqrt{3} + 6 = 0$, which factors as $\left(x - \sqrt{3}\right)\left(x - 2\sqrt{3}\right) = 0$.  Therefore, $x = \sqrt{3}$ and $x = 2\sqrt{3}$ are both possible values of $x$, so $x$ is $\fbox{(\textbf{E}) not uniquely determined}$.

 -j314andrews

 


	\item (\textit{AHSME 1980}) The number of teeth in three meshed circular gears $A$, $B$, and $C$ are $x$, $y$, and $z$, respectively. (The teeth on all gears are the same size and regularly spaced.) The angular speeds, in revolutions per minute, of $A$, $B$, and $C$ are in the proportion


 $\text{(A)} \ x: y: z ~~\text{(B)} \ z: y: x ~~ \text{(C)} \ y: z: x~~ \text{(D)} \ yz: xz: xy ~~ \text{(E)} \ xz: yx: zy$

 \textbf{Answer: }D



\textbf{Solution}

Since the gears are meshed, they must all rotate the same number of teeth per minute.  Let $k$ be this rate. Then gears $A$, $B$, and $C$ must have angular speeds of $\frac{k}{x}$, $\frac{k}{y}$, and $\frac{k}{z}$, respectively.  Thus the ratio of the three gears' angular speeds is $\frac{k}{x}:\frac{k}{y}:\frac{k}{z}$.  Multiplying each term of this ratio by $\frac{xyz}{k}$ yields $\boxed{(\textbf{D})\ yz:xz:xy}$.

 


	\item (\textit{AHSME 1980}) If the sum of the first $10$ terms and the sum of the first $100$ terms of a given arithmetic progression are $100$ and $10$, 
respectively, then the sum of first $110$ terms is


 $\text{(A)} \ 90 \qquad \text{(B)} \ -90 \qquad \text{(C)} \ 110 \qquad \text{(D)} \ -110 \qquad \text{(E)} \ -100$

 \textbf{Answer: }D



\textbf{Solution 1}

Let $a$ and $d$ be the first term and common difference of the sequence, respectively.

 The sum of the first $10$ terms is $\frac{10}{2}(2a+9d)=100$. This equation can be simplified to $2a+9d=20$.

 The sum of the first $100$ terms is $\frac{100}{2}(2a+99d)=10$. This equation can be simplified to $2a+99d=\frac{1}{5}$.

 Solving the system of these two equations yields $(a, d) = \left(\frac{1099}{100}, -\frac{11}{50}\right)$.  Therefore, the sum of the first $110$ terms is $\frac{110}{2}(2a+109d)=55\left(2\cdot\frac{1099}{100} + 109 \cdot -\frac{11}{50}\right) = 55 \cdot -2 = \boxed{(\textbf{D}) -110}$.

 



\textbf{Solution 2}

Let $a_1$, $a_2$, ..., $a_n$ be the terms of this arithmetic sequence, and let $d$ be its common difference.  Recall that in an arithmetic sequence with an even number of terms, the mean of the middle two terms equals the mean of all terms in the sequence.  

 The mean of the first $10$ terms is $\frac{a_5 + a_6}{2} = \frac{100}{10} = 10$, while the mean of the first $100$ terms is $\frac{a_{50}+a_{51}}{2} = \frac{10}{100} = \frac{1}{10}$.  

 Then $\frac{1}{10} = \frac{a_{50}+a_{51}}{2} = \frac{a_5 + 45d + a_6 + 45d}{2} = \frac{a_5+a_6}{2} + 45d = 10 + 45d$, so $d = -\frac{11}{50}$.  So the mean of the first $110$ terms is $\frac{a_{55}+a_{56}}{2} = \frac{a_{50}+5d+a_{51}+5d}{2} = \frac{a_{50}+a_{51}}{2}+5d=\frac{1}{10}+5\cdot-\frac{11}{50}=-1$.  Therefore, the sum of the first $110$ terms is $110 \cdot -1 =  \boxed{(\textbf{D}) -110}$.

 -j314andrews

 


	\item (\textit{AHSME 1980}) The equations of $L_1$ and $L_2$ are $y=mx$ and $y=nx$, respectively. Suppose $L_1$ makes twice as large of an angle with the horizontal (measured counterclockwise from the positive x-axis ) as does $L_2$, and that $L_1$ has 4 times the slope of $L_2$. If $L_1$ is not horizontal, then $mn$ is


 $\text{(A)} \ \frac{\sqrt{2}}{2} \qquad \text{(B)} \ -\frac{\sqrt{2}}{2} \qquad \text{(C)} \ 2 \qquad \text{(D)} \ -2 \qquad \text{(E)} \ \text{not uniquely determined}$

 \textbf{Answer: }C



\textbf{Solution 1 (Trigonometry)}

Since $L_1$ has $4$ times the slope of $L_2$, $m = 4n$.  Let $O = (0,0)$, $A = (1,0)$, $B = (1,n)$, $C = (1,m)$, and $\theta = \angle AOB = \angle BOC$.  Then $m = \tan 2\theta$ and $n = \tan \theta$.  

 Since $m = 4n$, $\tan 2\theta = 4\tan \theta$. By the tangent double-angle formula, $\dfrac{2\tan \theta}{1-\tan^2 \theta} = 4\tan \theta$. Cross-multiplying and collecting terms on one side yields $4\tan^3 \theta - 2\tan\theta = 0$, which factors as $2\tan \theta(2\tan^2\theta - 1) = 0$. Substituting $\tan\theta = n$ yields $2n(2n^2-1)=0$.

 Since line $L_1$ is not horizontal, $n \neq 0$.  So $2n^2 - 1 = 0$, and thus $n^2 = \frac{1}{2}$.  Therefore, $mn = 4n^2 = 4 \cdot \frac{1}{2} = \boxed{(\textbf{C}) \ 2}$.

 -e\_power\_pi\_times\_i, edited by j314andrews

 



\textbf{Solution 2 (Angle Bisector Theorem)}

Since $L_1$ has $4$ times the slope of $L_2$, $m = 4n$.  Let $O = (0,0)$, $A = (1,0)$, $B = (1,n)$, $C = (1,m) = (1, 4n)$.  By the Angle Bisector Theorem, $\frac{OC}{OA} = \frac{BC}{AB}$.  That is, $OC = \frac{3n}{n} = 3$.

 By the Pythagorean Theorem, $AC = \sqrt{3^2-1^2} = 2\sqrt{2}$.  So $m = 2\sqrt{2}$ and $n = \frac{2\sqrt{2}}{4} = \frac{\sqrt{2}}{2}$.  Therefore, $mn = \left(2\sqrt{2}\right)\left(\frac{\sqrt{2}}{2}\right) = \boxed{(\textbf{C}) \ 2}$.

 -j314andrews

 


	\item (\textit{AHSME 1980}) A bug (of negligible size) starts at the origin on the coordinate plane. First, it moves $1$ unit right to $(1,0)$. Then it makes a $90^\circ$ counterclockwise and travels $\frac 12$ a unit to $\left(1, \frac 12 \right)$. If it continues in this fashion, each time making a $90^\circ$ degree turn counterclockwise and traveling half as far as in the previous move, to which of the following points will it come closest? 


 $\text{(A)} \ \left(\frac 23, \frac 23 \right) \qquad \text{(B)} \ \left( \frac 45, \frac 25 \right) \qquad \text{(C)} \ \left( \frac 23, \frac 45 \right) \qquad \text{(D)} \ \left(\frac 23, \frac 13 \right) \qquad \text{(E)} \ \left(\frac 25, \frac 45 \right)$

 \textbf{Answer: }B



\textbf{Solution 1}

Each step the bug takes is half as long as its previous step.  So each horizontal step is $\frac{1}{4}$ as long as and in the opposite direction of the previous horizontal step.  Similarly, each vertical step is $\frac{1}{4}$ as long as and in the opposite direction the previous vertical step.

 Therefore, the $x$-coordinate is the sum of an infinite geometric series with first term $1$ and common ratio $-\frac{1}{4}$, so it must be $\frac{1}{1-\left(-\frac{1}{4}\right)} = \frac{4}{5}$.  Similarly, the $y$-coordinate is the sum of an infinite geometric series with first term $\frac{1}{2}$ and common ratio $-\frac{1}{4}$, so it must be $\frac{\frac{1}{2}}{1-\left(-\frac{1}{4}\right)}  = \frac{2}{5}$.  Therefore, the bug's path approaches $\boxed{(\textbf{B}) \left(\frac{4}{5}, \frac{2}{5}\right)}$




\textbf{Solution 2 (Complex Plane)}

Consider the bug's path as an infinite sequence of vectors in the complex plane.  Then each step can be represented by a complex number, with the first step being $1$ and the second step being $\frac{1}{2}i$. 

 Since each step is half as long as and rotated $90^{\circ}$ counterclockwise from the previous step, the complex number representing a step must be $\frac{1}{2}i$ times the complex number representing the previous step.

 Therefore, the bug's path can be represented by an infinite geometric series with first term $1$ and common ratio $\frac{1}{2}i$.  The sum of this series is $\frac{1}{1-\frac{1}{2}i} = \frac{2}{2-i} = \frac{2(2+i)}{(2-i)(2+i)} = \frac{4}{5} + \frac{2}{5}i$, which corresponds to $\boxed{(\textbf{B}) \left(\frac{4}{5}, \frac{2}{5}\right)}$

 -jaspersun, edited by j314andrews

 


	\item (\textit{AHSME 1980}) If the function $f$ defined by 

 \[f(x)=\frac{cx}{2x+3} ,\quad x\neq -\frac{3}{2} ,\] 
satisfies $x=f(f(x))$ for all real numbers $x$ except $-\frac{3}{2}$, then $c$ is


 $\text{(A)} \ -3 \qquad  \text{(B)} \ - \frac{3}{2} \qquad  \text{(C)} \ \frac{3}{2} \qquad  \text{(D)} \ 3 \qquad  \text{(E)} \ \text{not uniquely determined}$

 \textbf{Answer: }A



\textbf{Solution 1}

Since $f(x)=\frac{cx}{2x+3}$, $f(f(x)) = \frac{c^2x}{2cx+6x+9} = x$ for all $x \neq \frac{3}{2}$. 

 Multiplying both sides by $\frac{2cx+6x+9}{x}$ yields $c^2=2cx+6x+9$. That is, $x(2c+6) + (3+c)(3-c) = 0$. Therefore, both $2c+6 = 0$ and $(3+c)(3-c) = 0$, so $c= \boxed{(\textbf{A}) -3}$.

 



\textbf{Solution 2}

Since $f(x)=\frac{cx}{2x+3}$, $f(f(x)) = \frac{c^2x}{2cx+6x+9} = x$.  However, $f$ is undefined at $x=-\frac{3}{2}$. 

 Substituting $x=-\frac{3}{2}$ into $\frac{c^2x}{2cx+6x+9} = x$ yields $\frac{c}{2} = -\frac{3}{2}$, so $c= \boxed{(\textbf{A}) -3}$.

 



\textbf{Solution 3}

Since $f(x) = \frac{cx}{2x+3} = \frac{\frac{1}{2}cx}{x+\frac{3}{2}} = \frac{\frac{1}{2}cx + \frac{3}{4}c - \frac{3}{4}c}{x+\frac{3}{2}} = \frac{\frac{1}{2}c\left(x+\frac{3}{2}\right) - \frac{3}{4}c}{x+\frac{3}{2}} = \frac{1}{2}c - \frac{\frac{3}{4}c}{x+\frac{3}{2}}$, the graph of $f(x)$ must be a hyperbola with asymptotes $x=-\frac{3}{2}$ and $y=\frac{1}{2}{c}$.  

 Since $f(f(x)) = x$, $f(x) = f^{-1}(x)$ and the graph of $f(x)$ is symmetrical across the line $y=x$.  The reflections of the asymptotes of $f(x)$ across $y=x$ are $y = -\frac{3}{2}$ and $x=\frac{1}{2}c$, so $-\frac{3}{2} = \frac{1}{2}c$ and thus $c= \boxed{(\textbf{A}) -3}$.

 -j314andrews

 



\textbf{Solution 4}

We are given a function $f(x) = \frac{cx}{2x + 3}$ for $x \neq -\frac{3}{2}$, and it satisfies the functional equation $x = f(f(x))$ for all real numbers $x \neq -\frac{3}{2}$. We are tasked with finding the value of $c$.

 \textbf{Step 1: Calculate $f(f(x))$}

 We begin by calculating $f(f(x))$, which is the composition of the function $f(x)$ with itself. To do this, we substitute $f(x) = \frac{cx}{2x + 3}$ into itself:

 
 \[f(f(x)) = f\left( \frac{cx}{2x + 3} \right).\]

 Substitute $\frac{cx}{2x + 3}$ into the formula for $f$:

 
 \[f\left( \frac{cx}{2x + 3} \right) = \frac{c \left( \frac{cx}{2x + 3} \right)}{2 \left( \frac{cx}{2x + 3} \right) + 3}.\]

 Simplify the numerator:

 
 \[\text{Numerator} = c \times \frac{cx}{2x + 3} = \frac{c^2 x}{2x + 3}.\]

 Now simplify the denominator:

 
 \[\text{Denominator} = 2 \times \frac{cx}{2x + 3} + 3 = \frac{2cx}{2x + 3} + 3.\]

 To combine the terms in the denominator, express $3$ with a denominator of $2x + 3$:

 
 \[\frac{2cx}{2x + 3} + 3 = \frac{2cx + 3(2x + 3)}{2x + 3} = \frac{2cx + 6x + 9}{2x + 3} = \frac{(2c + 6)x + 9}{2x + 3}.\]

 Thus, we have:

 
 \[f(f(x)) = \frac{\frac{c^2 x}{2x + 3}}{\frac{(2c + 6)x + 9}{2x + 3}} = \frac{c^2 x}{(2c + 6)x + 9}.\]

 \textbf{Step 2: Set up the functional equation}

 We are given that $f(f(x)) = x$ for all $x \neq -\frac{3}{2}$. Therefore, we set the expression for $f(f(x))$ equal to $x$:

 
 \[\frac{c^2 x}{(2c + 6)x + 9} = x.\]

 \textbf{Step 3: Solve the equation}

 To eliminate the fraction, multiply both sides of the equation by $(2c + 6)x + 9$:

 
 \[c^2 x = x \left( (2c + 6)x + 9 \right).\]

 Expand both sides:

 
 \[c^2 x = (2c + 6)x^2 + 9x.\]

 Now, move all terms to one side of the equation:

 
 \[0 = (2c + 6)x^2 + 9x - c^2 x.\]

 Factor out $x$:

 
 \[0 = x \left( (2c + 6)x + 9 - c^2 \right).\]

 Since this equation must hold for all $x \neq 0$, the expression in parentheses must be equal to zero:

 
 \[(2c + 6)x + 9 - c^2 = 0.\]

 This simplifies to:

 
 \[(2c + 6)x + (9 - c^2) = 0.\]

 For this to hold for all $x \neq 0$, the coefficient of $x$ must be zero, and the constant term must also be zero. Thus, we have the system of equations:

 
            \begin{itemize}
        \item $2c + 6 = 0$
        \item $9 - c^2 = 0$
            \end{itemize}
            \textbf{Step 4: Solve for $c$}

 From $2c + 6 = 0$, we solve for $c$:

 
 \[2c = -6 \quad \Rightarrow \quad c = -3.\]

 Substitute $c = -3$ into $9 - c^2 = 0$:

 
 \[9 - (-3)^2 = 9 - 9 = 0.\]

 So, $c = -3$ satisfies both equations.

 \textbf{Final Answer:}

 The value of $c$ is $\boxed{\textbf{(A) -3}}$.

 


	\item (\textit{AHSME 1980}) A store prices an item in dollars and cents so that when $4\%$ sales tax is added, no rounding is necessary because the result is exactly $n$ dollars where $n$ is a positive integer. The smallest value of $n$ is


 $\text{(A)} \ 1 \qquad \text{(B)} \ 13 \qquad \text{(C)} \ 25 \qquad \text{(D)} \ 26 \qquad \text{(E)} \ 100$

 \textbf{Answer: }B



\textbf{Solution}

Let $x$ be the price of the item in cents.  The sales tax is $.04x = \frac{x}{25}$ cents, so $n=\frac{1}{100}\left( x+\frac{x}{25}\right)=\frac{26x}{2500}=\frac{13x}{1250}$. 

 Since $x$ is positive integer, the smallest possible integer value of $n$ occurs when $x=1250$, and in this case $n$ is $\fbox{\text{(\textbf{B}) 13}}$.

 


	\item (\textit{AHSME 1980}) Four of the eight vertices of a cube are the vertices of a regular tetrahedron. Find the ratio of the surface area of the cube to the surface area of the tetrahedron.


 
                    \begin{figure}[H]
                        \centering
                        \includegraphics[width=0.5\textwidth]{Problem/AHSME_I_1980_16_Problem_Figure_1.png}
                    \end{figure}
                    
$\text{(A)} \ \sqrt 2 \qquad \text{(B)} \ \sqrt 3 \qquad \text{(C)} \ \sqrt{\frac{3}{2}} \qquad \text{(D)} \ \frac{2}{\sqrt{3}} \qquad \text{(E)} \ 2$

 \textbf{Answer: }B



\textbf{Solution}

Without loss of generality, suppose the cube has side length $1$. By the Pythagorean Theorem, the side length of the tetrahedron is $\sqrt2$.  Therefore, each face of the tetrahedron has area $\frac{\left(\sqrt{2}\right)^2 \cdot \sqrt{3}}{4} = \frac{\sqrt{3}}{2}$ and its total surface area is $4\cdot\frac{\sqrt3}{2}=2\sqrt3$. The surface area of the cube is $6\cdot1^2=6$, so the ratio of the surface area of the cube to the surface area of the tetrahedron is $\frac{6}{2\sqrt3}=\boxed{(\textbf{B})\ \sqrt3}$.

 -aopspandy, edited by j314andrews

 


	\item (\textit{AHSME 1980}) Given that $i^2=-1$, for how many integers $n$ is $(n+i)^4$ an integer?


 $\text{(A)} \ \text{none} \qquad \text{(B)} \ 1 \qquad \text{(C)} \ 2 \qquad \text{(D)} \ 3 \qquad \text{(E)} \ 4$

 \textbf{Answer: }D



\textbf{Solution 1}

Expanding $(n+i)^4$ yields $n^4+4in^3-6n^2-4in+1 = (n^4-6n^2+1) + (4n^3-4n)i$.  This quantity is an integer if and only if $4n^3 - 4n = 4n(n+1)(n-1) = 0$, that is, if $n \in \{0, 1, -1\}$.  Therefore, there are $\boxed{(\textbf{D})\ 3}$ such values of $n$.

 -aopspandy, edited by j314andrews

 



\textbf{Solution 2}

For $(n+i)^4$ to be a real number, $n+i$ must be a scalar multiple of an eighth root of unity, that is, $n+i = r\left(\cos \frac{k\pi}{4} + i \sin \frac{k\pi}{4}\right)$ where $r \geq 0$ is a real number and $k$ is an integer such that $0 \leq k \leq 7$.  So $r \sin \frac{k\pi}{4} = 1$ and $\sin \frac{k\pi}{4}$ is positive.  Therefore, $k \in \{1, 2, 3\}$. 

 If $k = 1$, $r = \frac{1}{\sin\frac{\pi}{4}} = \sqrt{2}$, $n = \sqrt{2}\cos\frac{\pi}{4} = 1$ and $(n+i)^4 = (1+i)^4 = -4$.

 If $k = 2$, $r = \frac{1}{\sin\frac{\pi}{2}} = 1$, $n = \cos\frac{\pi}{2} = 0$ and $(n+i)^4 = i^4 = 1$.  

 If $k = 3$, $r = \frac{1}{\sin\frac{3\pi}{4}} = \sqrt{2}$, $n = \sqrt{2}\cos\frac{3\pi}{4} = -1$ and $(n+i)^4 = (-1+i)^4 = -4$.   

 In all these cases, $(n+i)^4$ is an integer, so $\{1, 0, -1\}$ are the $\boxed{(\mathbf{D})\ 3}$ possible values of $n$.

 


	\item (\textit{AHSME 1980}) If $b>1$, $\sin x>0$, $\cos x>0$, and $\log_b \sin x = a$, then $\log_b \cos x$ equals


 $\text{(A)} \ 2\log_b(1-b^{a/2}) ~~\text{(B)} \ \sqrt{1-a^2} ~~\text{(C)} \ b^{a^2} ~~\text{(D)} \ \frac 12 \log_b(1-b^{2a}) ~~\text{(E)} \ \text{none of these}$

 \textbf{Answer: }D



\textbf{Solution}

Since $\log_b \sin x = a$, $b^a=\sin x$.

 Let $c = \log_b \cos x$. Then $b^c=\cos x$.

 Since $\sin^2x+\cos^2x=1$, 
 \[(b^c)^2+(b^a)^2=1\]
 \[b^{2c}+b^{2a}=1\]
 \[b^{2c}=1-b^{2a}\]
 \[\log_b (1-b^{2a}) = 2c\]
 \[c=\boxed{(\textbf{D}) \ \frac 12 \log_b(1-b^{2a})}\]

 -aopspandy

 


	\item (\textit{AHSME 1980}) Let $C_1, C_2$ and $C_3$ be three parallel chords of a circle on the same side of the center. 
The distance between $C_1$ and $C_2$ is the same as the distance between $C_2$ and $C_3$. 
The lengths of the chords are $20, 16$, and $8$. The radius of the circle is 


 $\text{(A)} \ 12 \qquad  \text{(B)} \ 4\sqrt{7} \qquad  \text{(C)} \ \frac{5\sqrt{65}}{3} \qquad  \text{(D)}\ \frac{5\sqrt{22}}{2}\qquad \text{(E)}\ \text{not uniquely determined}$

 \textbf{Answer: }D



\textbf{Solution}

Place this circle in the coordinate plane such that the center is at the origin all three chords are to the right of, and parallel to, the $y$-axis.  The equation of this circle is $x^2 + y^2 = r^2$, where $r > 0$ is the radius of the circle. 

 Let $A$, $B$, and $C$ be the upper ends of chords $C_1$, $C_2$, and $C_3$, respectively.

 Since the $x$-axis is perpendicular to the three chords, it must bisect each chord.  Let $a$ be the $x$-coordinate $A$, and $d$ be the distance between $C_1$ and $C_2$, which is also the distance between $C_2$ and $C_3$.  Then $A$ is at $(a, 10)$, $B$ is at $(a+d, 8)$, and $C$ is at $(a+2d, 4)$.

 Since $A$, $B$, and $C$ all lie on this circle, the following equations must hold:
 \[\\a^2 + 100 = r^2\ (\textrm{i})\] 
 \[\\(a+d)^2 + 64 = r^2\ (\textrm{ii})\] 
 \[\\(a+2d)^2 + 16 = r^2\ (\textrm{iii})\]Subtracting equation $(\textrm{i})$ from equation $(\textrm{ii})$ yields:
 \[\\(a+d)^2 - a^2 - 36 = 0\]
 \[\\2ad + d^2 = 36\ (\textrm{iv})\]Subtracting equation $(\textrm{i})$ from equation $(\textrm{iii})$ yields:
 \[\\(a+2d)^2 - a^2 - 84 = 0\]
 \[\\4ad + 4d^2 = 84\]
 \[\\ad + d^2 = 21\ (\textrm{v})\]Subtracting $(\textrm{v})$ from $(\textrm{iv})$ yields:
 \[\\ad = 15\ (\textrm{vi})\]Subtracting $(\textrm{vi})$ from $(\textrm{v})$ yields $d^2 - 6 = 0$, so $d = \sqrt{6}$.

 Substituting $d = \sqrt{6}$ into $(\textrm{vi})$ yields $\sqrt{6}a = 15$, so $a = \frac{15}{\sqrt{6}} = \frac{5\sqrt{6}}{2}$.

 Substituting $a = \frac{5\sqrt{6}}{2}$ yields  $r^2 = \left(\frac{5\sqrt{6}}{2}\right)^2 + 100 = \frac{275}{2}$, so $r = \sqrt{\frac{275}{2}} = \boxed{(\textbf{D})\ \frac{5\sqrt{22}}{2}}$



	\item (\textit{AHSME 1980}) A box contains $2$ pennies, $4$ nickels, and $6$ dimes. Six coins are drawn without replacement, 
with each coin having an equal probability of being chosen. What is the probability that the value of coins drawn is at least $50$ cents? 


 $\text{(A)} \ \frac{37}{924} \qquad  \text{(B)} \ \frac{91}{924} \qquad  \text{(C)} \ \frac{127}{924} \qquad  \text{(D)}\ \frac{132}{924}\qquad \text{(E)}\ \text{none of these}$

 \textbf{Answer: }C



\textbf{Solution 1}

Each of the $\binom{12}{6} = 924$ possible combinations of $6$ coins is equally likely to be selected.

 If $6$ dimes are selected, the total will be $60$ cents.  There is only $\binom{6}{6} = 1$ way to select $6$ dimes.

 If $5$ dimes are selected, the total will be more than $50$ cents.  There are $\binom{6}{5} = 6$ ways to select the dimes, and $\binom{6}{1} = 6$ ways to select the other coin, for a total of $6 \cdot 6 = 36$ ways to select $5$ dimes and $1$ other coin.

 If $4$ dimes are selected, the total will be at least $50$ cents if and only if the other $2$ coins are both nickels.  There are $\binom{6}{4} = 15$ ways to select $4$ dimes and $\binom{4}{2} = 6$ ways to select $2$ nickels, for a total of $15 \cdot 6 = 90$ ways to select $4$ dimes and $2$ nickels.

 If $3$ or fewer dimes are selected, the total will be less than $50$ cents.

 Therefore, there are $1+36+90=127$ ways to select $6$ coins with a total of at least $50$ cents, so the probability is $\boxed{(\textbf{C})\ \frac{127}{924}}$.

 -j314andrews

 



\textbf{Solution 2}

We want the number of Successful Outcomes over the number of Total Outcomes. We want to calculate the total outcomes first. Since we have $12$ coins and we need to choose $6$, we have $\binom{12}{6}$ = $924$ Total outcomes. For our successful outcomes, we can have $(1) 1$ penny and $5$ dimes, $2$ nickels and $4$ dimes, $1$ nickel and $5$ dimes, or $6$ dimes. 

 For the case of $1$ penny and $5$ dimes, there are $\binom{6}{5}$ ways to choose the dimes and $2$ ways to choose the pennies. That is $6 \cdot 2 = 12$ successful outcomes. For the case of $2$ nickels and $4$ dimes, we have $\binom{6}{4}$ ways to choose the dimes and $\binom{4}{2}$ ways to choose the nickels. We have $15 \cdot 6$ = $90$ successful outcomes. For the case of $1$ nickel and $5$ dimes, we have $\binom{4}{1} \cdot \binom{6}{5} = 24$. Lastly, we have $6$ dimes and $0$ nickels, and $0$ pennies, so we only have one case. Therefore, we have $\dfrac {12 + 90 + 24 + 1}{924} = \dfrac{127}{924}$ = $\boxed{C}$



	\item (\textit{AHSME 1980}) 
                    \begin{figure}[H]
                        \centering
                        \includegraphics[width=0.5\textwidth]{Problem/AHSME_I_1980_21_Problem_Figure_1.png}
                    \end{figure}
                    


 In triangle $ABC$, $\measuredangle CBA=72^\circ$, $E$ is the midpoint of side $AC$, 
and $D$ is a point on side $BC$ such that $2BD=DC$; $AD$ and $BE$ intersect at $F$. 
The ratio of the area of $\triangle BDF$ to the area of quadrilateral $FDCE$ is


 $\text{(A)} \ \frac 15 \qquad  \text{(B)} \ \frac 14 \qquad  \text{(C)} \ \frac 13 \qquad  \text{(D)}\ \frac{2}{5}\qquad \text{(E)}\ \text{none of these}$

 \textbf{Answer: }A



\textbf{Solution}


                    \begin{figure}[H]
                        \centering
                        \includegraphics[width=0.5\textwidth]{Solution/AHSME_I_1980_21_Solution_Figure_1.png}
                    \end{figure}
                    

 Let $M$ be the midpoint of $\overline{DC}$.  Then $\triangle ECM \sim \triangle ACD$ and $\overline{EM} \parallel \overline{AD}$.  Since $\overline{EM} \parallel \overline{FD}$, it follows that $\triangle BFD \sim \triangle BEM$.  

 Let $a$ be the area of $\triangle BFD$.  Since the sides of $\triangle BEM$ are twice as long as the corresponding sides of $\triangle BFD$, the area of $\triangle BEM$ must be $2^2=4$ times the area of $\triangle BFD$, that is, $4a$.  

 Since the height of $\triangle BEC$ is the same as the height of $\triangle BEM$ and the base of $\triangle BEC$ is $\frac{3}{2}$ times the base of $\triangle BEM$, the area of $\triangle BEC$ is $\frac{3}{2}$ times the area of $\triangle BEM$, or $\frac{3}{2} \cdot 4a = 6a$.  

 Thus the area of quadrilateral $FDCE$ is $6a - a = 5a$, so the ratio of the area of $\triangle BFD$ to the area of quadrilateral $FDCE$ is $\frac{a}{5a} = \boxed{(\textbf{A})\ \frac{1}{5}}$.

 -j314andrews

 


	\item (\textit{AHSME 1980}) For each real number $x$, let $f(x)$ be the minimum of the numbers $4x+1, x+2$, and $-2x+4$. Then the maximum value of $f(x)$ is 


 $\text{(A)} \ \frac{1}{3} \qquad  \text{(B)} \ \frac{1}{2} \qquad  \text{(C)} \ \frac{2}{3} \qquad  \text{(D)} \ \frac{5}{2} \qquad  \text{(E)}\ \frac{8}{3}$

 \textbf{Answer: }E



\textbf{Solution}

Let $a(x) = 4x+1$, $b(x) = x+2$, and $c(x) = -2x+4$.  Then $f(x) = \min(a(x), b(x), c(x))$.

 $a(x) - b(x) = 3x - 1$, so $a(x) > b(x)$ when $x > \frac{1}{3}$, $a(x) = b(x)$ when $x = \frac{1}{3}$, and $a(x) < b(x)$ when $x < \frac{1}{3}$.

 $a(x) - c(x) = 6x - 3$, so $a(x) > c(x)$ when $x > \frac{1}{2}$, $a(x) = c(x)$ when $x = \frac{1}{2}$, and $a(x) < c(x)$ when $x < \frac{1}{2}$.

 $b(x) - c(x) = 3x - 2$, so $b(x) > c(x)$ when $x > \frac{2}{3}$, $b(x) = c(x)$ when $x = \frac{2}{3}$, and $b(x) < c(x)$ when $x < \frac{2}{3}$.

 If $x \leq \frac{1}{3}$, $a(x) \leq b(x) < c(x)$, so $f(x) = a(x) = 4x+1$.

 If $\frac{1}{3} \leq x \leq \frac{1}{2}$, $b(x) \leq a(x) \leq c(x)$, so $f(x) = b(x) = x+2$.

 If $\frac{1}{2} \leq x \leq \frac{2}{3}$, $b(x) \leq c(x) \leq a(x)$, so $f(x) = b(x) = x+2$.

 If $x \geq \frac{2}{3}$, $c(x) \leq b(x) < a(x)$, so $f(x) = c(x) = -2x+4$. 

 So $f(x)$ is increasing while $x < \frac{2}{3}$ and decreasing while $x > \frac{2}{3}$, and therefore, $f(x)$ is maximized at $x = \frac{2}{3}$, with a value of $f\left(\frac{2}{3}\right) = \frac{2}{3} + 2 = \boxed{(\textbf{E})\ \frac{8}{3}}$



	\item (\textit{AHSME 1980}) Line segments drawn from the vertex opposite the hypotenuse of a right triangle to the points trisecting the hypotenuse have lengths $\sin x$ and $\cos x$, where $x$ is a real number such that $0<x<\frac{\pi}{2}$. The length of the hypotenuse is 


 $\text{(A)} \ \frac{4}{3} \qquad  \text{(B)} \ \frac{3}{2} \qquad  \text{(C)} \ \frac{3\sqrt{5}}{5} \qquad  \text{(D)}\ \frac{2\sqrt{5}}{3}\qquad \text{(E)}\ \text{not uniquely determined}$

 \textbf{Answer: }C



\textbf{Solution 1}

Let $A$, $B$, and $C$ be the vertices of the right triangle, where $C$ is the right angle.  Let $a$, $b$, and $c$ be the lengths of the sides opposite $A$, $B$, and $C$, respectively.  Place $\triangle ABC$ in the coordinate plane such that $C$ is at the origin, $B$ lies on the positive $x$-axis, and $A$ lies on the positive $y$-axis.  Then $A$ has coordinates $(0, b)$ and $B$ has coordinates $(a, 0)$. 

 Let $M$ and $N$ be the trisection points of $\overline{AB}$, with $M$ closer to $A$ and $N$ closer to $B$.  Then $M$ has coordinates $\left(\frac{1}{3}a, \frac{2}{3}b\right)$ and $N$ has coordinates $\left(\frac{2}{3}a, \frac{1}{3}b\right)$.

 Let $m = CM$ and $n = CN$.  Then $m^2 = \frac{1}{9}a^2 + \frac{4}{9}b^2$ and $n^2 = \frac{4}{9}a^2 + \frac{1}{9}b^2$.  Adding these two equations yields $m^2 + n^2 = \frac{5}{9}(a^2+b^2) = \frac{5}{9}c^2$.

 Either $m = \sin x$ and $n = \cos x$, or $m = \cos x$ and $n = \sin x$.  In both cases, $m^2 + n^2 = \sin^2 x + \cos^2 x = 1$.  Therefore, $\frac{5}{9}c^2 = 1$ and $c = \sqrt{\frac{9}{5}} = \boxed{(\textbf{C})\ \frac{3\sqrt{5}}{5}}$.

 -j314andrews

 



\textbf{Solution 2 (Stewart's Theorem)}

Let $A$, $B$, and $C$ be the vertices of the right triangle, where $C$ is the right angle.  Let $a$, $b$, and $c$ be the lengths of the sides opposite $A$, $B$, and $C$, respectively. Let $M$ and $N$ be the trisection points of $\overline{AB}$.  Let $m = CM$ and $n = CN$.

 Using Stewart's Theorem on $\triangle ABC$ with segment $CM$ yields:

 
 \[m^2 \cdot c + \frac{2}{3}c \cdot \frac{1}{3}c \cdot c = a^2 \cdot \frac{1}{3}c + b^2 \cdot \frac{2}{3}c\]

 Dividing both sides by $c$ and isolating $m^2$ yields:

 
 \[m^2 = \frac{1}{3}a^2 + \frac{2}{3}b^2 - \frac{2}{9}c^2\ (\textrm{i})\]

 Using Stewart's Theorem on $\triangle ABC$ with segment $CN$ yields:

 
 \[n^2 \cdot c + \frac{2}{3}c \cdot \frac{1}{3}c \cdot c =a^2 \cdot \frac{2}{3}c + b^2 \cdot \frac{1}{3}c\]

 Dividing both sides by $c$ and isolating $n^2$ yields:

 
 \[n^2 = \frac{2}{3}a^2 + \frac{1}{3}b^2 - \frac{2}{9}c^2\ (\textrm{ii})\]

 Adding equations $(\textrm{i})$ and $(\textrm{ii})$ yields:

 
 \[m^2 + n^2 = a^2 + b^2 - \frac{4}{9}c^2\].

 Either $m = \sin x$ and $n = \cos x$, or $m = \cos x$ and $n = \sin x$.  In both cases, $m^2 + n^2 = \sin^2 x + \cos^2 x = 1$.  Also, by the Pythagorean theorem, $a^2 + b^2 = c^2$.

 Therefore, $1=c^2 - \frac{4}{9}c^2 = \frac{5}{9}c^2$, so $c^2 = \frac{9}{5}$ and $c = \sqrt{\frac{9}{5}} = \boxed{(\textbf{C})\ \frac{3\sqrt{5}}{5}}$.

 


	\item (\textit{AHSME 1980}) For some real number $r$, the polynomial $8x^3-4x^2-42x+45$ is divisible by $(x-r)^2$. Which of the following numbers is closest to $r$? 


 $\text{(A)} \ 1.22 \qquad  \text{(B)} \ 1.32 \qquad  \text{(C)} \ 1.42 \qquad  \text{(D)} \ 1.52 \qquad  \text{(E)} \ 1.62$

 \textbf{Answer: }D



\textbf{Solution 1}

Since $(x-r)^2$ is a factor of $8x^3-4x^2-42x+45$, $r$ is a double root. Let $s$ be the third root of $8x^3-4x^2-42x+45$. 

 By Vieta's formulas, $2r+s = \frac{1}{2}$, so $s = \frac{1}{2} - 2r$.  

 Also, by Vieta's formulas, $r^2 + 2rs = -\frac{21}{4}$.  Substituting $s = \frac{1}{2} - 2r$ yields $r^2 + 2r\left(\frac{1}{2} - 2r\right) = -3r^2 + r = -\frac{21}{4}$.  So $12r^2 - 4r - 21 = 0$, that is, $(6r + 7)(2r - 3) = 0$.  Therefore, $(r, s)$ is either $\left(\frac{3}{2}, -\frac{5}{2}\right)$ or $\left(-\frac{7}{6}, \frac{17}{6}\right)$. 

 Finally, by Vieta's formulas, $r^2s = -\frac{45}{8}$.  Though $(r, s) = \left(-\frac{7}{6}, \frac{17}{6}\right)$ does not satisfy this equation, $\left(r, s\right) = \left(\frac{3}{2}, -\frac{5}{2}\right)$ does.  

 Therefore, $r = \frac{3}{2} = 1.5$, so $\boxed{\text{(\textbf{D})} \ 1.52}$ is closest to $r$.

 -e\_power\_pi\_times\_i, edited by j314andrews

 



\textbf{Solution 2}

By polynomial long division, $\frac{8x^3 - 4x^2 - 42x + 45}{x^2 - 2rx + r^2} = 8x + (16r-4) + \frac{(24r^2 - 8r - 42)x + (-16r^3 + 4r^2 + 45)}{x^2 - 2rx + r^2}$.

 So $(24r^2 - 8r - 42)x + (-16r^3 + 4r^2 + 45) = 0$, that is, $r$ is a root of both $24r^2 - 8r - 42$ and $-16r^3 + 4r^2 + 45$. 

 Since $24r^2 - 8r - 42 = 2(6r + 7)(2r - 3)$, either $r = -\frac{7}{6}$ or $r = \frac{3}{2}$.  By the Rational Root Theorem, $-\frac{7}{6}$ is not a root of $-16r^3 + 4r^2 + 45$.  However, $-16\left(\frac{3}{2}\right)^3 + 4\left(\frac{3}{2}\right)^2 + 45 = -16 \cdot \frac{27}{8} + 4 \cdot \frac{9}{4} + 45 = -54 + 9 + 45 = 0$, so $\frac{3}{2}$ is a root of $-16r^3 + 4r^2 + 45$.

 Therefore, $r = \frac{3}{2} = 1.5$, so $\boxed{(\textbf{D})\ 1.52}$ is closest to $r$.

 -j314andrews

 



\textbf{Solution 3}

Let $y = 2x$.  Then $x = \frac{y}{2}$, so $8x^3 - 4x^2 - 42x + 45 = y^3 - y^2 - 21y + 45$.

 By the Rational Root Theorem, any rational root of $y^3 - y^2 - 21y + 45$ is a factor of $45$.  Of the factors of $45$, $3$ and $5$ are roots, and $y^3 - y^2 - 21y + 45 = (y-3)^2(y-5)$.  So $8x^3 - 4x^2 - 42x + 45 = (2x-3)^2(2x-5) = 8\left(x-\frac{3}{2}\right)^2\left(x-\frac{5}{2}\right)$, and $r = \frac{3}{2} = 1.5$. So $\boxed{(\textbf{D})\ 1.52}$ is closest to $r$.

 -j314andrews

 


	\item (\textit{AHSME 1980}) In the non-decreasing sequence of odd integers $\{a_1,a_2,a_3,\ldots \}=\{1,3,3,3,5,5,5,5,5,\ldots \}$ each odd positive integer $k$
appears $k$ times. It is a fact that there are integers $b, c$, and $d$ such that for all positive integers $n$, 
$a_n=b\lfloor \sqrt{n+c} \rfloor +d$, 
where $\lfloor x \rfloor$ denotes the largest integer not exceeding $x$. The sum $b+c+d$ equals 


 $\text{(A)} \ 0 \qquad  \text{(B)} \ 1 \qquad  \text{(C)} \ 2 \qquad  \text{(D)} \ 3 \qquad  \text{(E)} \ 4$

 \textbf{Answer: }C



\textbf{Solution 1}

Define the sequence $\{x_n\} = \lfloor\sqrt{n}\rfloor$ for all non-negative integers $n$.  Then for any non-negative integer $m$, $x_n = m$ if and only if $m \leq \sqrt{n} < m+1$, that is, $m^2 \leq n < (m+1)^2$.  So each non-negative integer $m$ appears $(m+1)^2 - m^2 = 2m+1$ times.  So $\{x_0, x_1, x_2, x_3, ...\} = \{0, 1, 1, 1, 2, 2, 2, 2, 2, ...\}$.  

 In the sequence $a_n$, each positive odd integer $2m+1$ appears $2m+1$ times.  Notice that $\{2x_0 + 1, 2x_1 + 1, 2x_2 + 1, 2x_3 + 1, ...\} = \{1, 3, 3, 3, 5, 5, 5, 5, 5, ...\} = \{a_1, a_2, a_3, a_4, ...\}$.  

 Therefore, $a_n = 2x_{n-1} + 1 = 2\lfloor{\sqrt{n-1}}\rfloor + 1$, so $b = 2$, $c = -1$, and $d = 1$, and $b + c + d = \boxed{(\textbf{C})\ 2}$.

 -j314andrews

 



\textbf{Solution 2}

Because the set consists of odd numbers, and since $\lfloor{}\sqrt{n+c}\rfloor{}$ is an integer and can be odd or even, $b = 2$ and $|a| = 1$. However, given that $\lfloor{}\sqrt{n+c}\rfloor{}$ can be $0$, $a = 1$. Then, $a_1 = 1 = 2\lfloor{}\sqrt{1+c}\rfloor{}+1$, and $\lfloor{}\sqrt{1+c}\rfloor{}$ = 0, and $c = -1$ because $c$ is an integer. $b+c+d = 2+(-1)+1 = \boxed{\text{(C)}\  2}$

 -e\_power\_pi\_times\_i

 


	\item (\textit{AHSME 1980}) Four balls of radius $1$ are mutually tangent, three resting on the floor and the fourth resting on the others. 
A tetrahedron, each of whose edges have length $s$, is circumscribed around the balls. Then $s$ equals 


 $\text{(A)} \ 4\sqrt 2 \qquad  \text{(B)} \ 4\sqrt 3 \qquad  \text{(C)} \ 2\sqrt 6 \qquad  \text{(D)}\ 1+2\sqrt 6\qquad \text{(E)}\ 2+2\sqrt 6$

 \textbf{Answer: }E



\textbf{Solution}

Let $A$, $B$, $C$, and $D$ be the centers of the four spheres of radius $1$.  These points must form a regular tetrahedron of side length $2$.  Let $O$ be the center of $\triangle ABC$. Then $OA = \frac{2\sqrt{3}}{3}$.  By the Pythagorean Theorem, the height $DO$ of tetrahedron $ABCD$ must be $\sqrt{2^2 - \left(\frac{2\sqrt{3}}{3}\right)^2} = \frac{2\sqrt{6}}{3}$.  Let $I$ be the center of tetrahedron $ABCD$.  Then $IABC$, $IABD$, $IACD$, and $IBCD$ are all congruent tetrahedra, each with $\frac{1}{4}$ of the volume of $ABCD$.  Since tetrahedra $IABC$ and $ABCD$ share base $ABC$, the height $IO$ of tetrahedron $IABC$, and therefore the inradius of tetrahedron $ABCD$, must be $\frac{1}{4} \cdot \frac{2\sqrt{6}}{3} = \frac{\sqrt{6}}{6}$.  

 Let $EFGH$ be the tetrahedron circumscribed around the spheres, with $AE$, $BF$, $CG$, and $DH$ intersecting at $I$.  Planes $ABC$ and $EFG$ must be parallel and a distance $1$ from each other.  So the distance from $I$ to plane $EFG$, and therefore the inradius of tetrahedron $EFGH$, must be $1 + \frac{\sqrt{6}}{6}$.  Therefore, the ratio of the inradius of tetrahedron $EFGH$ to the inradius of tetrahedron $ABCD$ is $\frac{1+\sqrt{6}/6}{\sqrt{6}/6} = 1 + \sqrt{6}$.  The ratio of the side length of tetrahedron $EFGH$ to the side length of tetrahedron $ABCD$ must also be $1+\sqrt{6}$, so $s = 2 (1+\sqrt{6}) = 2 + 2\sqrt{6}$. $\fbox{(E)}$

 -j314andrews 

 


	\item (\textit{AHSME 1980}) The sum $\sqrt[3] {5+2\sqrt{13}}+\sqrt[3]{5-2\sqrt{13}}$ equals 


 $\text{(A)} \ \frac 32 \qquad  \text{(B)} \ \frac{\sqrt[3]{65}}{4} \qquad  \text{(C)} \ \frac{1+\sqrt[6]{13}}{2} \qquad  \text{(D)}\ \sqrt[3]{2}\qquad \text{(E)}\ \text{none of these}$

 \textbf{Answer: }E



\textbf{Solution}

Let $y = \sqrt[3] {5+2\sqrt{13}}$, $z = \sqrt[3]{5-2\sqrt{13}}$, and $x = \sqrt[3] {5+2\sqrt{13}}+\sqrt[3]{5-2\sqrt{13}} = y + z$. 

 Then $x^3 = (y+z)^3 = y^3 + 3y^2z + 3yz^2 + z^3 = y^3 + z^3 + 3yz(y+z) = y^3 + z^3 + 3yzx$.

 Note that $y^3 + z^3 = 5 + 2\sqrt{13} + 5 - 2\sqrt{13} = 10$ and $yz = \sqrt[3]{\left(5 + 2\sqrt{13}\right)\left(5-2\sqrt{13}\right)} = \sqrt[3]{5^2 -\left(2\sqrt{13}\right)^2} = \sqrt[3]{-27} = -3$.

 So $x^3 = 10 + 3 \cdot -3 \cdot x = 10-9x$, that is $x^3+9x-10=0$.  This factors to $(x-1)(x^2+x+10) = 0$, which has $x=1$ as its only real solution.

 Therefore, $\sqrt[3] {5+2\sqrt{13}}+\sqrt[3]{5-2\sqrt{13}} = 1$ and the answer is $\boxed{(\textbf{E})\ \textrm{none of these}}$.

 


	\item (\textit{AHSME 1980}) The polynomial $x^{2n}+1+(x+1)^{2n}$ is not divisible by $x^2+x+1$ if $n$ equals 


 $\text{(A)} \ 17 \qquad  \text{(B)} \ 20 \qquad  \text{(C)} \ 21 \qquad  \text{(D)} \ 64 \qquad  \text{(E)} \ 65$

 \textbf{Answer: }C



\textbf{Solution 1}

The roots of $x^2+x+1$ are $-\frac{1}{2} + \frac{\sqrt{3}}{2}i = \cos\frac{2}{3}\pi + i \sin \frac{2}{3}\pi$ and $-\frac{1}{2} - \frac{\sqrt{3}}{2}i = \cos \frac{4}{3}\pi + i \sin\frac{4}{3}\pi$, which are the primitive third roots of unity.

 Let $w = -\frac{1}{2} + \frac{\sqrt{3}}{2}i = \cos\frac{2}{3}\pi + i \sin \frac{2}{3}\pi$.  Note that $-\frac{1}{2} - \frac{\sqrt{3}}{2}i = \cos \frac{4}{3}\pi + i \sin\frac{4}{3}\pi = w^2 = \bar{w}$.  So by the conjugate root theorem, $x^{2n} + 1 + (1+x)^{2n}$ is divisible by $x^2+x+1$ if and only if $w$ is a root of $x^{2n} + 1 + (1+x)^{2n}$.

 Also, $w^{2n} + 1 + (w+1)^{2n} = w^{2n} + 1 + (-w^2)^{2n} = w^{2n} + 1 + w^{4n} = w^{2n} + 1 + w^n = w^{2n} + w^n + 1$.

 If $n \equiv 0\ (\textrm{mod}\ 3)$, $w^{2n} + w^n + 1 \equiv 1 + 1 + 1 = 3$.

 If $n \equiv 1\ (\textrm{mod}\ 3)$, $w^{2n} + w^n + 1 \equiv w^2 + w + 1 = 0$.

 If $n \equiv 2\ (\textrm{mod}\ 3)$, $w^{2n} + w^n + 1 \equiv w + w^2 + 1 = 0$.

 Therefore, $x^{2n} + 1 + (1+x)^{2n}$ is not divisible by $x^2 + x + 1$ if and only if $n$ is divisible by $3$.  Of the answer choices, only $\boxed{(\textbf{C})\  21}$ is divisible by $3$.

 



\textbf{Solution 2}

Let $h(x)=x^2+x+1$.

 Then $(x+1)^{2n} = (x^2+2x+1)^n = (x+h(x))^n = \sum_{k=0}^{n} \binom{n}{k} x^k (h(x))^{n-k}$ $= \sum_{k=0}^{n-1} \binom{n}{k} x^k (h(x))^{n-k} + x^n = h(x) \left(\sum_{k=0}^{n-1} \binom{n}{k} x^k (h(x))^{n-k-1}\right) + x^n$.

 Let $g(x) = \sum_{k=0}^{n-1} \binom{n}{k} x^k (h(x))^{n-k-1}$.  Then $x^{2n} + 1 + (x+1)^{2n} = x^{2n} + 1 + g(x)h(x) + x^n$, so $x^{2n} + 1 + (x+1)^{2n}$ is divisible by $x^2 + x + 1$ if and only if $x^{2n} +  x^n + 1$ is divisible by $x^2 + x + 1$.

 Let $k$ be any non-negative integer, and let $r \in \{0, 1, 2\}$.  Then $x^{3k+r} - x^r = x^r(x^{3k} - 1) = x^r(x^3 - 1)\left(\sum_{j=0}^{k - 1} x^{3j}\right)$$= x^r(x-1)(x^2 + x + 1)\left(\sum_{j=0}^{k - 1} x^{3j}\right)$.

 Therefore, $x^{3k+r} - x^r$ is divisible by $x^2 + x + 1$.

 If $n$ is a nonnegative integer, then $n$ can be written as $3k + r$, where nonnegative integer $k$ and $r \in \{0, 1, 2\}$ are the quotient and remainder when $n$ is divided by $3$.  

 If $r = 2$, then $x^{2n} + x^n + 1 = x^{6k+4} + x^{3k+2} + 1 = x^{3(2k+1)+1} + x^{3k+2} + 1 = (x^{3(2k+1)+1} - x) + (x^{3k+2} - x^2) + x^2 + x + 1$.  Since $x^{3(2k+1)+1} - x$ and $(x^{3k+2} - x^2)$ are divisible by $x^2 + x + 1$, $x^{2n} + x^n + 1$ is divisible by $x^2 + x + 1$.  

 If $r = 1$, then $x^{2n} + x^n + 1 = x^{6k+2} + x^{3k+1} + 1 = x^{3(2k)+2} + x^{3k+1} + 1 = (x^{3(2k)+2} - x^2) + (x^{3k+1} - x) + x^2 + x + 1$.  Since $x^{3(2k)+2} - x^2$ and $(x^{3k+1} - x)$ are divisible by $x^2 + x + 1$, $x^{2n} + x^n + 1$ is divisible by $x^2 + x + 1$. 

 If $r = 0$, then $x^{2n} + x^n + 1 = x^{6k} + x^{3k} + 1 = (x^{6k} - 1) + (x^{3k} - 1) + 3$.  Since $x^{6k} - 1$ and $x^{3k} - 1$ are divisible by $x^2 + x + 1$, $x^{2n} + x^n + 1$ is not divisible by $x^2 + x + 1$.

 So $x^{2n} + 1 + (1+x)^{2n}$ is not divisible by $x^2 + x + 1$ if and only if $r = 0$, that is, $n$ is divisible by $3$.  The only answer choice that is divisible by $3$ is $\boxed{(\textbf{C})\ 21}$.

 -Wei, edited by j314andrews

 



\textbf{Solution 3}

We start by noting that 
 \[x + 1 \equiv -x^2 \mod (x^2+x+1).\]Let $n = 3k+r$, where $r \in \{ 0,1,2 \}$.

 Thus we have 
 \[x^{4n} + x^{2n} + 1 \equiv x^{4r} + x^{2r} + 1 \mod (x^3 -1).\]

 When $r = 0$, 
 \[x^{4n} + x^{2n} + 1 \equiv 3 \mod (x^3 -1).\]When $r = 1$, 
 \[x^{4n} + x^{2n} + 1 \equiv x^2 + x + 1 \mod (x^3 -1),\] which will be divisible by $x^2+x+1$.

 When $r = 2$, 
 \[x^{4n} + x^{2n} + 1 \equiv x^2 + x + 1 \mod (x^3 -1),\] which will also be divisible by $x^2+x+1$.

 Thus $r \ne 0$, so $n$ cannot be divisible by $3$, and the answer is $\textrm{(C)}$.

 


	\item (\textit{AHSME 1980}) How many ordered triples $(x,y,z)$ of integers satisfy the system of equations below? 


 
 \[\begin{array}{l} x^2-3xy+2y^2-z^2=31 \\ -x^2+6yz+2z^2=44 \\ x^2+xy+8z^2=100\\ \end{array}\]
$\text{(A)} \ 0 \qquad  \text{(B)} \ 1 \qquad  \text{(C)} \ 2 \qquad  \text{(D)}\ \text{a finite number greater than 2}\qquad \text{(E)}\ \text{infinitely many}$

 \textbf{Answer: }A



\textbf{Solution 1 (Sum of Squares)}

We are considering the sum of three equations:

 
 \[x^2 - 2xy + 2y^2 + 6yz + 9z^2 = (x - y)^2 + (y + 3z)^2 = 175.\]

 Since $x, y, z$ are integers, this implies that $175$ can be written as the sum of two squares, i.e.,

 
 \[175 = a^2 + b^2,\]

 where $a$ and $b$ are integers. 

 However, $175 \equiv 3\ (\mathrm{mod}\ 4)$, but $a^2$ and $b^2$ must each be congruent to either $0$ or $1 (\mathrm{mod}\ 4)$. So this equation, and therefore this system, has $0$ integer solutions. $\fbox{(A)}$

 \href{https://artofproblemsolving.com/wiki/index.php?title=User:Wwei.yu&action=edit&redlink=1}{Wwei.yu} (\href{https://artofproblemsolving.com/wiki/index.php?title=User\_talk:Wwei.yu&action=edit&redlink=1}{talk}) 22:09, 28 March 2020 (EDT)

 



\textbf{Solution 2 (Eliminate Cross Terms)}

First, to eliminate the $xy$ term, multiply both sides of the third equation by $3$ to get $3x^2 + 3xy + 24z^2 = 300$, and then add it to the first equation to get $4x^2 + 2yz + 23z^2 = 331$.  

 Next, to eliminate the $yz$ term, multiply both sides of the new equation by $-3$ to get $-12x^2 - 6yz - 69z^2 = -993$, and then add it to the second equation to get $-13x^2 - 67z^2 = -949$.  

 This equation can be rearranged to get $67z^2 = 949-13x^2 = 13(73-x^2)$.  Since $67$ is not divisible by $13$, $z$ must be divisible by $13$ and there exists an integer $n$ such that $z = 13n$.  Substituting $z = 13n$ and dividing both sides of this equation by $13$ yields $871n^2 = 73 - x^2$.  So $871n^2 \leq 73$, which is only true if $n = 0$.  But $n = 0$ would require $x^2 = 73$, which is impossible since $73$ is not a perfect square.  Therefore, this equation (and also this system) has $0$ integer solutions. $\fbox{(A)}$

 -j314andrews, based on solution by Farenhajt

 



\textbf{Solution 3 (Modular Arithmetic)}

Reducing the second equation mod $2$ yields $x^2 \equiv 0\ (\mathrm{mod}\ 2)$.  Therefore, $x \equiv 0\ (\mathrm{mod}\ 2)$ and $x$ is even.

 Reducing the first equation mod $2$ and substituting yields $z^2 \equiv 1\ (\mathrm{mod}\ 2)$.  Therefore, $z \equiv 1\ (\mathrm{mod}\ 2)$ and $z$ is odd.

 Since $x$ is even, $x^2 \equiv 0\ (\mathrm{mod}\ 4)$, and since $z$ is odd, $z^2 \equiv 1\ (\mathrm{mod}\ 4)$. Also $6z \equiv 2z \equiv 2\ (\mathrm{mod}\ 4)$. Reducing the second equation mod $4$ and substituting yields $2y + 2 \equiv 0\ (\mathrm{mod}\ 4)$, so $y$ is odd.

 Reducing the third equation mod $4$ and substituting yields $xy \equiv 0\ (\mathrm{mod}\ 4)$.  Since $y$ is odd, $x \equiv 0\ (\mathrm{mod}\ 4)$.  

 Since $y$ is odd, $y^2 \equiv 1\ (\mathrm{mod}\ 4)$.  Reducing the first equation mod $4$ and substituting yields $1 \equiv 3\ (\mathrm{mod}\ 4)$, which is impossible.

 Therefore, this system has $0$ integer solutions. $\fbox{(A)}$

 -j314andrews

 


	\item (\textit{AHSME 1980}) A six digit number (base $10$) is squarish if it satisfies the following conditions: 


 (i) none of its digits are zero; 


 (ii) it is a perfect square; and 


 (iii) the first of two digits, the middle two digits and the last two digits of the number are all perfect squares when considered as two digit numbers. 


 How many squarish numbers are there? 


 $\text{(A)} \ 0 \qquad  \text{(B)} \ 2 \qquad  \text{(C)} \ 3 \qquad  \text{(D)} \ 8 \qquad  \text{(E)} \ 9$

 \textbf{Answer: }B



\textbf{Solution}

If $N$ is a squarish number, then $N = 10000a^2 + 100b^2+ c^2$, where $a$, $b$, and $c$ are all integers between $4$ and $9$ inclusive.  The only way that $N$ could be a perfect square would be if $N = (100a + c)^2$.  Expansion yields $N = 10000a^2 + 200ac + c^2$, so $2ac = b^2$, and therefore $b$ must be even.

 If $b = 4$, then $ac = 8$, which is impossible.

 If $b = 6$, then $ac = 18$, which is impossible.

 If $b = 8$, then $ac = 32$, which is possible if $(a,c)$ is either $(8,4)$ or $(4,8)$.  This yields squarish numbers $646416 = 804^2$ and $164664 = 408^2$.

 Therefore, there are $2$ squarish numbers. $\fbox{(B)}$ 

 \href{https://artofproblemsolving.com/wiki/index.php?title=User:Wwei.yu&action=edit&redlink=1}{Wwei.yu} (\href{https://artofproblemsolving.com/wiki/index.php?title=User\_talk:Wwei.yu&action=edit&redlink=1}{talk})Wei

 -Edited by j314andrews and wzheng05

 


\end{enumerate}